\appendix
\section{经典无失真 MVDR 的求解证明}\label{app:classic}
在经典 MVDR 问题中，固定目标接收导向矢量 $\bm a\in\mathbb C^{N\times1}$ 和干扰加噪声协方差矩阵 $\mathbf R$，假设 $\bm a\ne\bm0$、$\mathbf R=\mathbf R^H\succ0$。令接收向量为 $\bm v$，则经典无失真 MVDR 问题为
\begin{align}\label{eq:app-mvdr-problem}
\min_{\bm v\in\mathbb C^{N\times1}}\quad
&\bm v^H\mathbf R\bm v,\\
\text{s.t.}\quad &\bm v^H\bm a=1.\notag
\end{align}
其中目标函数是输出干扰与噪声功率，约束要求目标方向的复数响应为 $1$。这一问题及其闭式解可参见 Capon [1] 和 Vorobyov [2，第 2.2 节]。

\noindent\textbf{证明。}对于任意满足约束 $\bm v^H\bm a=1$ 的向量 $\bm v$，由于 $\mathbf R\succ0$，其正定平方根和逆平方根均存在。将单位矩阵 $\mathbf R^{1/2}\mathbf R^{-1/2}$ 插入内积，并使用 Cauchy--Schwarz 不等式，得到
\begin{align}\label{eq:app-classic-cs}
1&=\left|\bm v^H\bm a\right|^2
=\left|\bigl(\mathbf R^{1/2}\bm v\bigr)^H
\bigl(\mathbf R^{-1/2}\bm a\bigr)\right|^2\notag\\
&\le\bigl(\bm v^H\mathbf R\bm v\bigr)
\bigl(\bm a^H\mathbf R^{-1}\bm a\bigr).
\end{align}
由于 $\bm a\ne\bm0$ 且 $\mathbf R^{-1}\succ0$，有 $\bm a^H\mathbf R^{-1}\bm a>0$，所以任何可行向量都满足
\begin{equation}\label{eq:app-classic-bound}
\bm v^H\mathbf R\bm v\ge
\frac{1}{\bm a^H\mathbf R^{-1}\bm a}.
\end{equation}
令
\begin{equation}\label{eq:app-classic-solution}
\bm v^{\rm MVDR}=
\frac{\mathbf R^{-1}\bm a}{\bm a^H\mathbf R^{-1}\bm a}.
\end{equation}
由于 $\bm a^H\mathbf R^{-1}\bm a$ 为正实数，直接代入可得
\begin{align}\label{eq:app-classic-check}
\bigl(\bm v^{\rm MVDR}\bigr)^H\bm a&=1,\\
\bigl(\bm v^{\rm MVDR}\bigr)^H\mathbf R\bm v^{\rm MVDR}
&=\frac{1}{\bm a^H\mathbf R^{-1}\bm a}.\notag
\end{align}
因此，$\bm v^{\rm MVDR}$ 满足无失真约束并达到式~\eqref{eq:app-classic-bound} 的下界，是问题~\eqref{eq:app-mvdr-problem} 的全局最优解。由 Cauchy--Schwarz 取等条件及无失真约束，这一解唯一。

在本文的单位范数接收约束下，将式~\eqref{eq:app-classic-solution} 的方向归一化，得到
\begin{equation}\label{eq:app-unit-norm}
\bm g^{\rm unit}=\frac{\bm v^{\rm MVDR}}{\|\bm v^{\rm MVDR}\|}
=\frac{\mathbf R^{-1}\bm a}{\|\mathbf R^{-1}\bm a\|}.
\end{equation}
对于非零接收向量，$|\bm v^H\bm a|^2/(\bm v^H\mathbf R\bm v)$ 不受非零复数缩放影响；因此两种归一化对应相同的最大 SINR 方向。经典 MVDR 的 $\bm v^{\rm MVDR}$ 满足无失真约束，本文的 $\bm g^{\rm unit}$ 满足单位范数约束，它们的数值一般不相同。若采用原式分母中的常数噪声项，只有单位范数向量才可直接用 $\bm g^H\mathbf R\bm g$ 表示该分母。\hfill$\square$

\section{MVDR 成立条件及本文模型的对应关系}\label{app:conditions}
本节分别说明固定其他变量时的最大 SINR 结论，以及在 $P0$ 联合优化中可以采用 MVDR 接收波束的条件。后一结论还依赖感知发射波束的配合调整。

\noindent\textbf{一、最大感知 SINR 的条件及本文依据。}
\begin{enumerate}[label=\textbf{(\arabic*)},leftmargin=*]
\item \textbf{秩一目标回波与非零接收导向。}式~\eqref{eq:A} 将感知信道写为接收导向向量与发射导向向量的外积，因而固定其他变量时，感知发射波束只通过标量增益影响 SINR；论文给出的接收导向向量首元素为 $1$，故不为零。若发射目标增益为零，所有接收方向的 SINR 均为零，MVDR 仍是最大值的一个解，但不再是唯一的最大方向。
\item \textbf{干扰加噪声矩阵正定。}原问题中的匹配指标非负、CU 发射功率为正；进一步采用 $\sigma^2>0$ 的正噪声功率假设，则式~\eqref{eq:new-cov} 给出的矩阵为正定矩阵，逆矩阵存在。论文将 $\sigma^2$ 定义为噪声功率，但应用结论时仍需明确其严格为正。
\item \textbf{接收向量的定义域与归一化。}论文本身把波束定义为复向量，并约束 $\|\bm g_i(t)\|=1$；因此可使用式~\eqref{eq:mvdr} 的单位范数形式。原 $P0$ 变量行中三个波束的 $\ge0$ 与复向量定义不一致；根据作者已确认的笔误，应将其理解为复向量域，而不施加逐元素非负实数约束。
\item \textbf{信道与目标方向的可用性。}计算该闭式接收波束需要所用时隙的目标接收导向向量及干扰加噪声矩阵。论文明确假设目标相对方位角已知，信道在时隙内准静态；但准静态并不自动意味着干扰信道被准确估计。若要声称实际计算所得波束达到真实模型的最大 SINR，应另外说明相关信道或协方差和噪声功率已准确获得；有估计误差时，该精确最优性不能直接保证。
\end{enumerate}

\noindent\textbf{二、将 MVDR 用于 $P0$ 最优方案的附加条件及本文依据。}
\begin{enumerate}[label=\textbf{(\arabic*)},leftmargin=*]
\item \textbf{感知发射波束可以缩放。}论文列出的感知发射约束是功率上限，且无已列出的发射功率下限；第三节构造将波束乘以 $[0,1]$ 内的系数，故不会违反该上限。该论证以 $P0$ 已列出完整的波束约束为前提；若另有不能保持于缩放下的硬件或发射约束，需要重新论证。
\item \textbf{保持感知数据量并控制通信干扰。}原模型中的感知速率由感知 SINR 决定；第三节同时调整接收和发射波束，使 SINR 不变，故感知数据量、感知卸载与计算负担不变。感知发射波束缩小不会增大同频 CU 的第一阶段干扰，其余阶段速率和变量不变，因而 CU 的成功卸载约束不恶化。
\item \textbf{能耗权重和最优解存在性。}目标值不增的构造需要 $\omega_2\ge0$。原论文只明确给出权重之和为 $1$，该等式本身不能推出各权重非负；应在模型中补充非负权重假设。第三节还以 $P0$ 存在全局最优解为前提，所得结论是：若最优解存在，则至少有一个采用 MVDR 接收方向。现有证明并未单独证明 $P0$ 一定存在全局最优解。
\end{enumerate}

综上，本文的秩一感知信道、接收单位范数和所列功率约束支持上述推导；正噪声功率、所需信道信息可用、非负能耗权重及 $P0$ 最优解存在性等条件，应按上述区别明确作为假设或进一步验证。
